\textit{Characteristics: Content describes the subject discussed, while question nature describes the purpose of the inquiry. The examples cover all 19 content categories and all five nature categories. Labels illustrate assistant interpretation rather than independent validation. English text summarizes the Chinese wording. Chinese text preserves the source wording after removing HTML and normalizing whitespace. Longer replies are explicitly marked as excerpts. Firm names follow the supplier records. Examples were not selected using stock returns.}

\subsubsection*{Content categories}

\noindent \textbf{Example 10: Financial report} \\
\textbf{平安银行} [Stock Code: 000001]. Question: 2 June 2010. Reply: 7 June 2010. \\
\textit{Content: Financial report. Source record: 6.} \\
\textbf{Question:} Will lending quotas affect the company's performance in the second through fourth quarters? (贵公司受放款额度所限是否会影响2-4季度的业绩？) \\
\textbf{Reply (excerpt):} First quarter lending focused more on small and medium enterprises, trade finance, and personal credit. Lending to government financing platforms, project loans, and certain restricted industries was tightly controlled. (我行1季度的贷款发放的结构更加倾向对中小企业、贸易融资和个人信贷的支持，严格限制政府平台授信、项目贷款和两高行业等。) \\
\textbf{Classification rationale:} The question concerns expected financial performance. The response discusses the lending composition behind that assessment.

\vspace{0.4cm}
\noindent \textbf{Example 11: Regulation} \\
\textbf{海王生物} [Stock Code: 000078]. Question: 1 September 2010. Reply: 8 September 2010. \\
\textit{Content: Regulation. Source record: 464.} \\
\textbf{Question:} Last year, the company said that the influenza vaccine would receive a production permit between June and August. Why has it not received the permit? When will it be obtained? (去年说甲流疫苗今年6-8月份获得生产许可，为什么现在还没有获得？是否欺骗了股民？到底什么时候能获？) \\
\textbf{Reply (excerpt):} The vaccine has not entered production or commercial launch. The subsidiary adjusted its work schedule in response to the progress of its cooperation with GSK and developments in the influenza epidemic. (公司控股子公司海王英特龙的疫苗至今没有投产上市，主要是因为海王英特龙根据其自身与GSK合作的进程及甲流的疫情发展情况进行的工作进度调整。) \\
\textbf{Classification rationale:} The inquiry asks about a production permit and its timing. This distinguishes regulatory approval from a general product inquiry.

\clearpage
\noindent \textbf{Example 12: External governance} \\
\textbf{中国宝安} [Stock Code: 000009]. Question: 5 January 2010. Reply: 6 January 2010. \\
\textit{Content: External governance. Source record: 67.} \\
\textbf{Question:} Are fund companies among the ten largest holders of tradable shares? Do institutional investors often visit the company for research? (敢问贵公司前十大流通股东有没有基金公司？请问有没有机构投资者经常去公司实地调研?) \\
\textbf{Reply:} Please consult the third quarter report for the list of the ten largest holders of tradable shares and information about institutional research visits. (公司前十大流通股东名单请查看公司第三季度报告。 机构投资者到公司实地调研情况也请查看公司第三季度报告。) \\
\textbf{Classification rationale:} The inquiry concerns external investors and their monitoring activity.

\vspace{0.4cm}
\noindent \textbf{Example 13: Asset restructuring} \\
\textbf{中国宝安} [Stock Code: 000009]. Question: 25 January 2010. Reply: 3 February 2010. \\
\textit{Content: Asset restructuring. Source record: 69.} \\
\textbf{Question:} When will Zhongbao Holdings, controlled by China Baoan, inject assets into Shenzhen Hongji? (请问深宝安直接控股的中宝控股何时向深鸿基资产注入) \\
\textbf{Reply:} It is not convenient to reply. (不便回复) \\
\textbf{Classification rationale:} The question concerns an asset injection. The response declines to provide the requested timing.

\clearpage
\noindent \textbf{Example 14: Infringement, disputes, lawsuit} \\
\textbf{富奥股份} [Stock Code: 000030]. Question: 22 January 2010. Reply: 1 February 2010. \\
\textit{Content: Infringement, disputes, lawsuit. Source record: 154.} \\
\textbf{Question:} Why did the company not announce the Shenzhen Intermediate Court lawsuit on January 20? The company clearly knew about it. (为什么不公告1月20日深圳中院的官司？公司明明知道这件事情) \\
\textbf{Reply:} The amount involved is small and does not meet the disclosure threshold. It should be handled under the cumulative disclosure threshold in subsequent work. (经查询得知，该案涉及金额很小，未达披露标准，宜按累计披露标准在今后工作中处理。) \\
\textbf{Classification rationale:} The inquiry concerns a lawsuit and its disclosure, rather than a general regulatory rule.

\vspace{0.4cm}
\noindent \textbf{Example 15: Violation} \\
\textbf{东旭蓝天} [Stock Code: 000040]. Question: 4 November 2010. Reply: 16 November 2010. \\
\textit{Content: Violation. Source record: 288.} \\
\textbf{Question:} On November 4, the company received a notice of investigation from the Shenzhen inspection bureau of the securities regulator for suspected false statements. This is disappointing. (公司于11月4日收到证监会深圳稽查局《立案调查通知书》，因公司涉嫌虚假陈述，深圳稽查局决定对公司及相关人员立案调查。 真叫人失望啊。) \\
\textbf{Reply:} The company will cooperate with the inspection and disclose the results promptly. (公司将积极配合深圳稽查局的稽查工作，并对相关稽查结果及时履行信息披露义务。) \\
\textbf{Classification rationale:} The subject is an investigation of an alleged violation. The allegation is part of the source text, not a finding of this study.

\clearpage
\noindent \textbf{Example 16: Stock trading} \\
\textbf{万 科Ａ} [Stock Code: 000002]. Question: 9 May 2010. Reply: 20 May 2010. \\
\textit{Content: Stock trading. Source record: 30.} \\
\textbf{Question:} Apart from property regulation, are there other reasons for the recent decline in Vanke A shares? What measures has the company taken, and how should investors view its prospects? (最近万科A股票贬值除了房控的影响，还有别的原因吗？贵公司有采取什么措施，应怎么看待未来的万科？) \\
\textbf{Reply (excerpt):} Short term stock prices depend on macroeconomic and policy developments, capital market conditions, industry changes, funding conditions, and investor views. Property regulation is currently an important influence. (股票的价格短期表现受很多因素的影响，宏观经济和政策走势、资本市场的格局、行业的变化、资金面的变动情况、投资者的心态等等，目前房地产的调控政策是重要的影响因素。) \\
\textbf{Classification rationale:} The inquiry starts from the company's secondary market share price.

\vspace{0.4cm}
\noindent \textbf{Example 17: Insider trading} \\
\textbf{方大集团} [Stock Code: 000055]. Question: 24 November 2010. Reply: 25 November 2010. \\
\textit{Content: Insider trading. Source record: 376.} \\
\textbf{Question:} Can the company disclose the buyers and sellers in the block trade on November 23, 2010? Did it involve insider trading or price manipulation? (能否公布2010年11月23日大宗交易的买卖双方，是否存在内幕交易或操纵股价) \\
\textbf{Reply:} The company does not know the identities of the buyers and sellers and cannot disclose them. It states that the trade did not involve insider trading. (我公司不知晓2010年11月23日大宗交易的买卖双方的情况，无法公布买卖双方的情况，该笔大宗交易不存在内幕交易的情况。) \\
\textbf{Classification rationale:} The investor explicitly asks about insider trading and manipulation. The company's denial is not independently verified here.

\clearpage
\noindent \textbf{Example 18: Dividend} \\
\textbf{国农科技} [Stock Code: 000004]. Question: 13 January 2010. Reply: 18 January 2010. \\
\textit{Content: Dividend. Source record: 55.} \\
\textbf{Question:} The company has not paid dividends or distributed shares for ten years. Does it plan to do so this year? (公司已经有10年没有分红或送股了,是两市仅有铁公鸡,今年可有分红或送股了的安排?) \\
\textbf{Reply:} A forecast of a loss was published at the end of the third quarter. Whether the company will be profitable this year remains uncertain. Investors should invest rationally and consider the risks. (本公司三季度末已发布预亏公告，今年是否盈利暂时还不能确定，敬请投资者理性投资，注意风险。) \\
\textbf{Classification rationale:} The request concerns a dividend policy. The response relates the possibility of a distribution to profitability.

\vspace{0.4cm}
\noindent \textbf{Example 19: Financing} \\
\textbf{平安银行} [Stock Code: 000001]. Question: 15 March 2010. Reply: 26 March 2010. \\
\textit{Content: Financing. Source record: 3.} \\
\textbf{Question:} Shenzhen Development Bank A has paid a dividend only once in seven years. Does this make it ineligible to apply for a private placement? (深发展A（000001）七年才分了一次红,是否不能申请增发融资？) \\
\textbf{Reply (excerpt):} The bank signed a private placement agreement with Ping An Life and announced the proposal on June 13, 2009. It meets the relevant eligibility conditions. (我行已与平安人寿签订定向增发协议，并于2009年6月13日公告了关于向平安人寿非公开发行股票的方案，公司符合关于非公开发行股票的有关条件。) \\
\textbf{Classification rationale:} The main subject is a financing transaction. The dividend history supplies the premise of the inquiry.

\clearpage
\noindent \textbf{Example 20: Investment} \\
\textbf{世纪星源} [Stock Code: 000005]. Question: 5 August 2010. Reply: 20 August 2010. \\
\textit{Content: Investment. Source record: 60.} \\
\textbf{Question:} The company's carport project has been underway for more than ten years. Has it been completed? What progress has been made? (世纪星源“车港项目”搞了10多年，竣工没有?进展情况是怎样的？) \\
\textbf{Reply:} The company disclosed the progress on August 11. Please consult the listed company announcement on CNINFO. (关于车港项目的进展情况，本司于8月11日作了披露，详细内容请参阅巨潮网站的上市公司公告。) \\
\textbf{Classification rationale:} The inquiry concerns the completion of an investment project.

\vspace{0.4cm}
\noindent \textbf{Example 21: Operation-operating assets} \\
\textbf{中国宝安} [Stock Code: 000009]. Question: 18 November 2010. Reply: 22 November 2010. \\
\textit{Content: Operation-operating assets. Source record: 75.} \\
\textbf{Question:} How many grams of graphene can the subsidiary produce? How long will it take to reach mass production? (子公司石墨稀产量到底是多少克？量产要多长时间？) \\
\textbf{Reply:} Research is still underway. Mass production has not started, and its timing cannot be determined. (目前正在研发之中，尚未形成量产，具体达到量产的时间无法确定。) \\
\textbf{Classification rationale:} The quantity and production capacity are the focus. The case is close to the product development category, so this illustrative label depends on the emphasis of the inquiry.

\clearpage
\noindent \textbf{Example 22: Operation-products or business} \\
\textbf{中国宝安} [Stock Code: 000009]. Question: 23 December 2010. Reply: 23 December 2010. \\
\textit{Content: Operation-products or business. Source record: 77.} \\
\textbf{Question:} Has the subsidiary BTR successfully developed graphene products? When will mass production begin? (请问贵公司控股子公司贝特瑞新材料产品石墨烯是否已经研发成功？何时量产？谢谢) \\
\textbf{Reply:} Development is still in progress. The timing of mass production is undetermined. (尚在研发之中，量产时间待定。) \\
\textbf{Classification rationale:} The focus is whether a particular product has been developed. A shared topic does not imply a complete answer or positive sentiment.

\vspace{0.4cm}
\noindent \textbf{Example 23: Operation-industry related} \\
\textbf{中集集团} [Stock Code: 000039]. Question: 6 June 2010. Reply: 12 June 2010. \\
\textit{Content: Operation-industry related. Source record: 192.} \\
\textbf{Question:} CNOOC plans to purchase substantial offshore equipment. What is the basis of its cooperation with CIMC? What competitive position does CIMC hold in the domestic offshore market? (中海油今后几年将购置大量的海工设备,中集与中海油的合作基础如何?对于国内海工市场,中集处于什么样的竞争地位?谢谢!) \\
\textbf{Reply:} CIMC and CNOOC have begun cooperating. CIMC Raffles holds contracts for three semisubmersible platforms and expects to deliver one this year. Management describes its position in domestic offshore equipment as relatively strong. (目前中集与中海油已展开合作，中集烟台莱佛士已经持有中海油的3个半潜式钻井平台的建造合同，并预计在今年会交付其中的1台。在国内海工市场,中集烟台莱佛士在手持订单量、建造关键设备、详细设计、施工设计能力经验等方面，处于相对领先的竞争地位。) \\
\textbf{Classification rationale:} The inquiry asks about the firm's competitive position in an industry, not just one product.

\clearpage
\noindent \textbf{Example 24: Operation-macroeconomic environment} \\
\textbf{中集集团} [Stock Code: 000039]. Question: 25 May 2010. Reply: 26 May 2010. \\
\textit{Content: Operation-macroeconomic environment. Source record: 187.} \\
\textbf{Question:} Has the European debt crisis affected the company's business? Has euro depreciation caused direct losses, and has container demand been affected? (请问最近的欧洲债危机是否对公司业务有影响；欧元贬值是否造成直接经济损失；最近公司集装箱业务是否受到欧洲债务危机的冲击) \\
\textbf{Reply (excerpt):} Management expects effects on exports, receivables, exchange losses, and financing. It also expects benefits from importing European materials and says that financial instruments can hedge currency risk. (最近的欧洲债务危机预计将对本公司业务产生一定的影响：首先会影响中国对欧洲出口，人民币相对升值，出口成本增加，出口将可能减少；其次，本公司所涉及的欧洲业务其应收账款减值、汇兑损失可能增加，在企业融资方面也可能产生不利影响。 \ldots{} 不过从正面来看，本公司将可以从进口欧洲原材料中受益，从而降低生产成本。 另一方面，本公司可以通过金融工具规避和对冲欧元贬值的风险，所以不会造成很大的直接经济损失。) \\
\textbf{Classification rationale:} The question relates a macroeconomic shock to company operations. The reply describes both adverse and favorable exposures.

\vspace{0.4cm}
\noindent \textbf{Example 25: Operation-others} \\
\textbf{东旭蓝天} [Stock Code: 000040]. Question: 29 January 2010. Reply: 2 February 2010. \\
\textit{Content: Operation-others. Source record: 234.} \\
\textbf{Question:} The company's operating ability, development, and profitability appear weak. Does the chair have a plan to reverse this position in the first quarter? (公司的经营能力、发展能力和盈利水平如此的差,请问董事长今年一季度有计划扭转这种劣势吗?) \\
\textbf{Reply:} The company is preparing to divest businesses unrelated to property and become a specialized property company. Investors should consult the first quarter report for its operating performance. (公司目前进行相关准备工作，逐步实现非房地产业务的剥离，致力打造专业化房地产公司，提升盈利水平；公司2010年第一季度经营情况敬请留意《2010年第一季度报告》。) \\
\textbf{Classification rationale:} The inquiry asks about the overall operating plan. It does not isolate a specific product, asset, or macroeconomic shock.

\clearpage
\noindent \textbf{Example 26: Corporate governance} \\
\textbf{平安银行} [Stock Code: 000001]. Question: 16 November 2010. Reply: 19 November 2010. \\
\textit{Content: Corporate governance. Source record: 12.} \\
\textbf{Question:} I am a shareholder. Why am I not allowed to attend the shareholder meeting? (我是你们的股东，为什么不让参加股东大会？) \\
\textbf{Reply:} All shareholders registered on the record date may attend the meeting and may appoint a proxy in writing to attend and vote. The proxy need not be a shareholder. (根据我行股东大会通知，凡于该次股东大会股权登记日在中国证券登记结算有限责任公司深圳分公司登记在册的本公司全体股东均有权出席股东大会，并可以书面形式委托代理人出席会议和参加表决，该股东代理人不必是本公司股东。) \\
\textbf{Classification rationale:} The subject is shareholder participation in the firm's internal governance.

\vspace{0.4cm}
\noindent \textbf{Example 27: Stock repurchase} \\
\textbf{晨鸣纸业} [Stock Code: 000488]. Question: 28 June 2010. Reply: 29 June 2010. \\
\textit{Content: Stock repurchase. Source record: 878.} \\
\textbf{Question:} Last year's shareholder meeting approved an H share repurchase plan. Why has nothing happened after more than a year? (晨鸣纸业去年股东大会通过回购H股计划，为什么到现在已一年有余，仍不见动静？) \\
\textbf{Reply (excerpt):} The meeting on May 26, 2009 authorized the board to repurchase some H shares. The authorization lasted one year and expired automatically. (2009年5月26日召开的2008年度股东大会审议通过了授权董事会回购公司部分H股股份，授权为一年的期限，期限届满后自动失效，详情请查看公司2008年度股东大会决议公告。) \\
\textbf{Classification rationale:} The response explains an earlier repurchase authorization rather than announcing a new repurchase.

\clearpage
\noindent \textbf{Example 28: Others} \\
\textbf{国农科技} [Stock Code: 000004]. Question: 9 December 2010. Reply: 10 December 2010. \\
\textit{Content: Others. Source record: 57.} \\
\textbf{Question:} What is the company's website? The published address is no longer accessible. (请问贵公司的网址是什么，目前公布的网址已上不去了) \\
\textbf{Reply:} The company is applying for a new domain and rebuilding its website. Completion is expected by the end of the year. (正在申请新的域名，重新建设网站，预计年底完成。) \\
\textbf{Classification rationale:} Website access is a valid investor inquiry outside the other eighteen subjects.

\clearpage
\subsubsection*{Question nature categories}

\noindent \textbf{Example 29: Answers or explanations} \\
\textbf{平安银行} [Stock Code: 000001]. Question: 2 June 2010. Reply: 7 June 2010. \\
\textit{Question nature: Answers or explanations. Source record: 6.} \\
\textbf{Question:} Will lending quotas affect the company's performance in the second through fourth quarters? (贵公司受放款额度所限是否会影响2-4季度的业绩？) \\
\textbf{Reply (excerpt):} First quarter lending focused more on small and medium enterprises, trade finance, and personal credit. (我行1季度的贷款发放的结构更加倾向对中小企业、贸易融资和个人信贷的支持，严格限制政府平台授信、项目贷款和两高行业等。) \\
\textbf{Classification rationale:} The investor asks management to explain how an operating constraint affects performance.

\vspace{0.4cm}
\noindent \textbf{Example 30: Correct misunderstandings} \\
\textbf{TCL 集团} [Stock Code: 000100]. Question: 21 April 2010. Reply: 21 April 2010. \\
\textit{Question nature: Correct misunderstandings. Source record: 512.} \\
\textbf{Question:} Can shareholders profit from a private placement with a floor price of RMB 3.46? Why does the company not raise the floor to reflect the market price? (以3.46元为底价的定向增发，大股东们(就算第一大股东除外)可利用低价增发套利是不是事实？公司为何根据市场价不提高定向增发底价？) \\
\textbf{Reply:} RMB 3.46 is only the floor price. The final issue price is determined through investor inquiry or bidding. (3.46只是发行底价,最终发行价格是需要投资者询价(或竟价)确定的.) \\
\textbf{Classification rationale:} The reply distinguishes the floor price from the final issue price and corrects the premise that shares will necessarily be issued at the floor.

\clearpage
\noindent \textbf{Example 31: Verify rumors} \\
\textbf{东旭蓝天} [Stock Code: 000040]. Question: 26 February 2010. Reply: 2 March 2010. \\
\textit{Question nature: Verify rumors. Source record: 241.} \\
\textbf{Question:} A rumor says that China Baoan will inject property assets into Shenzhen Hongji A. Is this true? (有传闻称中国宝安即将向深鸿基A注入地产资产!是否属实?) \\
\textbf{Reply:} The board has not received a letter from the largest shareholder, China Baoan Group Holdings, concerning an asset injection. (截止目前，公司董事局尚未收到公司第一大股东中国宝安集团控股有限公司关于资产注入的函件。) \\
\textbf{Classification rationale:} The question asks management to verify a specific rumor. It is not a general request for a restructuring plan.

\vspace{0.4cm}
\noindent \textbf{Example 32: Inappropriate questions} \\
\textbf{南 玻Ａ} [Stock Code: 000012]. Question: 30 July 2010. Reply: 30 July 2010. \\
\textit{Question nature: Inappropriate questions. Source record: 100.} \\
\textbf{Question:} The recently announced investment project did not disclose its expected benefits. What is the payback period of the new ultrathin electronic glass project? (我们注意到近期公布的投资项目未披露有关投资效益的内容，请问，河北南玻新建超薄 电子玻璃项目的投资回收期是多久？) \\
\textbf{Reply (excerpt):} The project's payback period is commercially confidential and cannot be provided. (项目的投资回收期属于商业机密，不便回复。) \\
\textbf{Classification rationale:} The question requests a project metric that the inquiry describes as undisclosed and the reply treats as confidential. This label identifies the request, not a legal finding that it is improper.

\clearpage
\noindent \textbf{Example 33: Suggestions or others} \\
\textbf{方大集团} [Stock Code: 000055]. Question: 9 July 2010. Reply: 12 July 2010. \\
\textit{Question nature: Suggestions or others. Source record: 371.} \\
\textbf{Question:} I suggest listing the three main businesses separately. As a shareholder for ten years, I am disappointed by dividends and market returns and hope the company will take action. (建议公司把三项主业分拆上市，作为方大的老股东十年来，无论是分红还是二级市场回报均令人失望，希望公司有所作为，也对投资人有所交待。) \\
\textbf{Reply:} The company currently has no plan to list the three businesses separately. (公司目前没有将三项产业分拆上市的计划。) \\
\textbf{Classification rationale:} The investor proposes an action rather than primarily requesting an explanation of existing information.

